%------------------------------------------------------------------------------%
\section{Álgebra Linear}
\label{sec:Algebra_Linear}

Apresentamos nessa seção algumas definições e propriedades estudadas na Álgebra Linear.
A finalidade destes resultados é generalizar o conceito de vetores, da Geometria Analítica.
Em essência dizemos que qualquer entidade matemática que tenha propriedades semelhantes as
propriedades dos vetores da geometria também são vetores, com isso podemos utilizar
resultados oriundos da geometria em contextos mais abstratos.

Para a Álgebra Linear as propriedades fundamentais dos vetores são as operações que
podemos realizar com eles. Assim definimos um espaço vetorial como sendo um conjunto
de elementos nos quais podemos realizar operações semelhantes as operações com vetores.
O termo espaço é utilizado sempre que tivermos um conjunto onde é possível
realizar alguma operação em todos os seus elementos. A definição a seguir explicita
o que é um espaço vetorial.

%------------------------------------------------------------------------------%
\begin{definicao}{Espaço Vetorial}{def:espaco-vetorial}
  Um \emph{Espaço Vetorial}\index{Espaço Vetorial} real é um conjunto $V$,
  não vazio, onde as operações de \emph{Soma}\index{Soma} e
  \emph{Produto por Escalar}\index{Produto!por escalar} estão definidas.
  Além disso, para quaisquer $u$, $v$ e $w\in V$ e $\alpha$ e $\beta\in\R$
  é sempre verdade que
  \begin{enumerate}
    \item \( v + w \in V \)
    \item \( \alpha v \in V \)
    \item \( (u+v) + w = u + (v+w)\)
    \item \( v+w = w+v \)
    \item Existe \( 0\in V\) tal que \(v+0=v\)
    \item Existe \(-v\in V\) tal que \(v+(-v) = 0\)
    \item \( \alpha(v+w) = \alpha v + \alpha w \)
    \item \((\alpha+\beta)v = \alpha v + \beta v \)
    \item \((\alpha \beta)v = \alpha(\beta v)\)
    \item \( 1v = v \)
  \end{enumerate}
\end{definicao}
%------------------------------------------------------------------------------%

Não importa a natureza dos elementos ou das operações utilizadas,
se as propriedades forem atendidas dizemos que $V$ é um \emph{Espaço Vetorial},
$v\in V$ é um \emph{Vetor}\index{Vetor} e $\alpha\in\R$ é um
\emph{Escalar}\index{Escalar}. A definição de espaços vetoriais complexos é
equivalente substituindo os escalares reais por números complexos. No contexto dessa
apostila estamos interessados em espaços vetoriais onde os vetores são funções
definidas em um intervalo, que chamaremos de
\emph{Espaços de Funções}\index{Espaço de Funções}.

Uma característica importante dos espaços vetoriais, reais ou complexos,
é que novos vetores podem ser construídos pela combinação linear de vetores dados.

%------------------------------------------------------------------------------%
\begin{definicao}{Combinação Linear}{def:combinacao-linear}
  Dados os vetores \(v_1, v_2, \dots, v_n \in V\) e escalares
  \( \alpha_1, \alpha_2, \dots, \alpha_n \) o vetor
  \[
    v = \alpha_1 v_1 + \alpha_2 v_2 + \cdots + \alpha_n v_n
  \]
  também é um elemento de $V$ e é chamado de
  \emph{Combinação Linear}\index{Combinação Linear} dos vetores~$v_i$.
\end{definicao}
%------------------------------------------------------------------------------%

Uma questão que muitas vezes precisamos responder é se um vetor
é ou não gerado por uma combinação linear de outros vetores.
Se ele puder ser gerado dessa forma, dizemos que ele é
Linearmente Dependente dos vetores
que o geram e em geral ele pode ser considerado supérfluo. Se tivermos
um conjunto de vetores onde nenhum pode ser gerado pelos demais dizemos
que esse conjunto é linearmente independente, como descrito na definição a seguir.

%------------------------------------------------------------------------------%
\begin{definicao}{Independência Linear}{def:independencia-linear}
  Dados os vetores \(v_1, v_2, \dots, v_n \in V\), dizemos que o conjunto
  \(\{v_1, v_2, \dots, v_n\}\) é
  \emph{Linearmente Independente}\index{Linearmente!independente} (LI)
  se a equação
  \[
   \alpha_1 v_1 + \alpha_2 v_2 + \cdots + \alpha_n v_n = 0
  \]
  implica que \( \alpha_1 = \alpha_2 = \cdots = \alpha_n = 0 \).
  Se existir uma solução onde algum escalar não é nulo o conjunto de vetores
  é \emph{Linearmente Dependente}\index{Linearmente!dependente} (LD).
\end{definicao}
%------------------------------------------------------------------------------%

Uma consequência da propriedade de que vetores podem ser gerados pela combinação
de outros é que podemos selecionar um conjunto de vetores para gerar todo o
espaço vetorial. O menor conjunto que gera todo o espaço é chamado de
base, como descrito na próxima definição.

%------------------------------------------------------------------------------%
\begin{definicao}{Base de um Espaço Vetorial}{def:base}
  Um conjunto \(\{v_1, v_2, \dots, v_n\}\) de vetores em $V$ é uma
  \emph{Base}\index{Base} de $V$ se:
  \begin{enumerate}
    \item for um conjunto linearmente independente
    \item todo vetor de $V$ pode ser gerado como combinação linear dos
          vetores do conjunto.
  \end{enumerate}
\end{definicao}
%------------------------------------------------------------------------------%

Um tipo de operação que pode ser estudada no contexto da Álgebra Linear são
as transformações lineares, que podem ser definidas como segue.

%------------------------------------------------------------------------------%
\begin{definicao}{Transformação Linear}{def:transformacao-linear}
  Dados dois espaços vetoriais $V$ e $W$, uma
  \emph{Transformação Linear}\index{Transformação Linear} é uma função de
  $V$ em $W$, \(T:V\to W\) que satisfaz:
  \begin{enumerate}
    \item para quaisquer $v$ e $w\in V$
    \[
      T(v+w) = T(v) + T(w)
    \]
    \item Para qualquer escalar $\alpha$ e $v\in V$
    \[
      T(\alpha v) = \alpha T(v)
    \]
  \end{enumerate}
\end{definicao}
%------------------------------------------------------------------------------%

Derivadas e integrais são exemplos de transformações lineares no espaço de funções,
assim como, as equações diferenciais lineares. Um problema importante que
surge no contexto de transformações lineares é o problema de autovalores e autovetores.

%------------------------------------------------------------------------------%
\begin{definicao}{Autovalores e Autovetores}{def:autovalores-autovetores}
  Dada uma transformação linear de um espaço vetorial $V$ nele mesmo,
  \(T:V\to V\), se existir um \emph{vetor não nulo} $v\in V$
  e um escalar $\lambda$ tais que
  \[
    T(v) = \lambda v
  \]
  dizemos que $\lambda$ é um \emph{Autovalor}\index{Autovalor} de $T$ e
  $v$ é o \emph{Autovetor}\index{Autovetor} de $T$ associado a $\lambda$.
\end{definicao}
%------------------------------------------------------------------------------%

No contexto dos espaços de função é comum chamar os autovetores de
\emph{Autofunções}\index{Autofunção}.

Além das operações de soma e produto por escalar, a Álgebra Linear
também generaliza outros conceitos como comprimento e ângulo entre
vetores.
Nesse momento é útil fazermos uma comparação com o plano cartesiano, $\R^2$,
onde calculamos o produto escalar, ou produto interno,
dos vetores $v=(v_1, v_2)$ e $w=(w_1, w_2)$ como
\[
  v\cdot w = \prodint{v, w} = v_1 w_1 + v_2 w_2
\]
Com o produto escalar podemos calcular a norma, ou tamanho, de um vetor $v\in\R^2$
como
\[
  \norm{v} = \sqrt{\prodint{v, v}\,} = \sqrt{ v_1^2 + v_2^2\,}
\]
Com a norma podemos calcular a distância entre dois vetores, ou pontos,
\[
  d(v, w) = \norm{v-w} = \sqrt{ (v_1-w_1)^2 + (v_2-w_2)^2\,}
\]
Podemos também calcular o ângulo entre dois vetores usando a relação entre
o produto escalar e a norma
\[
  v\cdot w = \norm{v}\, \norm{w}\, \cos\theta
\]
Finalmente dizemos que os vetores $v$ e $w$ são ortogonais se o ângulo entre
eles for \ang{90}, ou seja, se
\[
  v\cdot w = 0
\]
Vamos agora generalizar essas definições para espaços vetoriais quaisquer.

%------------------------------------------------------------------------------%
\begin{definicao}{Produto Interno}{def:produto-interno}
  Dado um espaço vetorial real $V$, um \emph{Produto Interno}\index{Produto!interno}
  sobre $V$ é uma função, \(\prodint{v, w}\),
  que associa um número real a cada par de vetores $v$ e $w\in V$,
  satisfazendo as propriedades
  \begin{enumerate}
    \item \(\prodint{v, v} \geq 0\) para todo $v\in V$
    \item \(\prodint{v, v} = 0\) se, e somente se, $v=0$
    \item \(\prodint{\alpha v, w} = \alpha \prodint{v, w}\) para todo escalar $\alpha$
    \item \(\prodint{u+v, w} = \prodint{u, w} + \prodint{v, w}\)
    \item \(\prodint{v, w} = \prodint{w, v}\)
  \end{enumerate}
\end{definicao}
%------------------------------------------------------------------------------%

Se o espaço vetorial $V$ possui uma definição de produto interno dizemos que ele é um
\emph{Espaço Vetorial com Produto Interno}\index{Espaço Vetorial!com produto interno}.
No estudo da transformada de Fourier utilizamos o produto interno
\begin{equation}
  \label{eq:produto_interno_funcoes}
  \prodint{v, w} = \int_a^b v(x) \, w(x) \, dx
\end{equation}
para o espaço de funções no intervalo $[a, b]$

%------------------------------------------------------------------------------%
\begin{definicao}{Ortogonalidade}{def:ortogonalidade}
  Em um espaço vetorial com produto interno $\prodint{,}$, dizemos
  que dois vetores $v$ e $w$ são \emph{Ortogonais}\index{Ortogonalidade} se
  \[
    \prodint{v, w} = 0
  \]
\end{definicao}
%------------------------------------------------------------------------------%

Para indicar que os vetores $v$ e $w$ são ortogonais podemos utilizar a notação
\(v\perp w\). A proposição a seguir apresenta as propriedades de vetores
ortogonais.

%------------------------------------------------------------------------------%
\begin{proposicao}{Propriedades da Ortogonalidade}{pro:propriedades-ortogonalidade}
  Em um espaço vetorial com produto interno, $V$, valem as propriedades de ortogonalidade
  \begin{enumerate}
    \item \( 0 \perp v \) para todo \(v\in V\)
    \item \( v \perp w \) implica que \( w \perp v \)
    \item Se \( v \perp w \) para todo \(w\in V\), então \(v=0\)
    \item Se \( u \perp w \) e \( v \perp w \), então  \( (u+v) \perp w \)
    \item Se \( v \perp w \), então \( (\lambda v) \perp w \) para todo escalar $\lambda$
  \end{enumerate}
\end{proposicao}
%------------------------------------------------------------------------------%

No caso do espaço de funções a condição de ortogonalidade oriunda do produto
interno~\eqref{eq:produto_interno_funcoes} é
\[
  \prodint{v, w} = \int_a^b v(x) \, w(x) \, dx = 0
\]
No estudo das séries de Fourier usamos o fato de que as funções seno e
cosseno com a escolha adequada de frequências são ortogonais.

Se o espaço vetorial possui um produto interno podemos definir uma norma,
como descrito na próxima definição

%------------------------------------------------------------------------------%
\begin{definicao}{Norma}{def:norma}
  Dado um espaço vetorial com produto interno \(\prodint{\;,\;}\), definimos a
  \emph{Norma}\index{Norma} de um vetor $v\in V$ em relação a esse produto interno
  como
  \[
    \norm{v} = \sqrt{\prodint{v, v}}
  \]
\end{definicao}
%------------------------------------------------------------------------------%

Se \(\norm{v}=1\) para algum vetor $v$ dizemos que ele é um
\emph{Vetor Unitário}\index{Vetor!unitário}.

A norma induzida no espaço de funções pelo produto
interno~\eqref{eq:produto_interno_funcoes} é
\[
  \norm{v} = \prodint{v, v} = \int_a^b v(x)^2 \, dx = 0
\]

Podemos agora definir a distância entre dois vetores em um espaço vetorial.

%------------------------------------------------------------------------------%
\begin{definicao}{Distância}{def:distancia}
  Dado um espaço vetorial com produto interno \(\prodint{,}\), definimos a
  \emph{Distância}\index{Distância} entre vetores $v$ e $w\in V$
  em relação a esse produto interno como
  \[
    d(v, w) = \norm{v-w}
  \]
\end{definicao}
%------------------------------------------------------------------------------%

A distância induzida no espaço de funções pelo produto
interno~\eqref{eq:produto_interno_funcoes} é
\[
  d(v, w) = \norm{v-w} = \int_a^b \big(v(x)-w(x)\big)^2 \, dx = 0
\]
Essa definição de distância é usada na Seção~\ref{sec:fenomeno_gibbs},
sobre o fenômeno de Gibbs, para esclarecer o significado apropriado a convergência
da série de Fourier.

%------------------------------------------------------------------------------%
